% Transcription — fonds Alexandre Grothendieck, shelfmark 23, batch 3 (pages 41-60).
% Established from the facsimile archives/batches/23/batch-03.pdf (a mirror of
% https://grothendieck.umontpellier.fr/23.pdf), per the skill
% transcribe-grothendieck. The batch file carries no cover sheet and opens at
% page 41: PDF page N is archivists' page 40+N. Checked against the pencilled
% numbers bottom-left, which read 41, 42, 43 ... 60 on the twenty leaves in
% order — the alignment holds end to end.
%
% Pass: Fable 5.1 (claude-fable-5-1), 2026-09-21 — first pass, unchecked against the pages by a human.
%
% What the batch is. Four blocks, three of them not in his hand, sorted by
% the covers he slipped between them.
%
%   41-42  Two blank leaves (a blue separator sheet, then a cream one). Absent.
%   43-46  A letter from Lê Dũng Tráng, Paris, 11 January 1968 (page 43), with
%          three sheets of errata numbered ①②③ in his hand (44-46): corrections
%          to EGA IV §§20-21 — the sheaf S(O_X) of regular sections, the sheaf
%          of meromorphic functions K_X, the section s_D of O_X(D), and the
%          notation of 21.6.2. This is the "lettres (1968)" of the inventory.
%          The letter is transcribed with its errata: it dates and attributes
%          them, and says which kind of statement they touch.
%   47     Blank cream leaf. Absent.
%   48     One leaf in his hand, ink over a faint printed verso, headed
%          « EGA IV 18 » and then « EGA VI »: two statements about sections of
%          an unramified X over Y — gluing a section over a schematically dense
%          open with a section over a closed subscheme with the same space, and
%          the exactness of X(Y) → X(Y') ⇉ X(Y'') for a finite surjective
%          epimorphism Y' → Y — with a marginal note asking to drop the
%          noetherian hypothesis and replace « fini » by « entier ». The hand is
%          fast and several connective words are lost.
%   49     An otherwise blank leaf inscribed « Suites régulières / Th. 11.2.9 de
%          Raynaud » in his hand, top right: the cover of the run that follows.
%          No \page; its words open the section.
%   50-51  Two IHES letterhead sheets in his hand: page 50 lists the places in
%          chapter IV where 11.2.9 is applied; page 51 lists the paragraphs on
%          regular sequences and regular immersions, points to the two Deligne
%          letters of October 1966, and carries a boxed question — struck by two
%          diagonals — on regular immersions in the non-noetherian case.
%   52-53  A typescript carbon, two pages, a « PS » in his own voice (« Dans mon
%          séminaire sur Riemann-Roch cette année ... »): the definition of
%          regular immersion in EGA IV 16.9.2 is not the right one outside the
%          noetherian case; the serviceable notion is the Koszul one; what to
%          do about it on the proofs of IV 16 and 19. Complete on the two pages.
%          The addressee is not named; the text says only what it says.
%   54     An otherwise blank leaf inscribed « Lettres Deligne » in his hand,
%          top right: cover of the last block. No \page.
%   55-60  Three letters from Pierre Deligne, in Deligne's hand, filed in
%          reverse order: Bruxelles, 15 décembre 1966 (55-56, the second leaf a
%          short continuation); Rösrath, 29 octobre 1966 (57-58); Bruxelles,
%          17 octobre 1966 (59-60, Deligne's pages 1 and 2 — the letter breaks
%          off mid-proof at the foot of 60 and continues in batch 4). They
%          answer his reading of EGA IV 19 and 0_IV: separation hypotheses,
%          19.5.6, 0_IV 19.3.12, Auslander's formula dim proj M = prof A − prof M,
%          a Tor-completion proposition, 16.1.5. His own pencil annotations on
%          them (« Err », « oui », « Th. de Auslander », « dans réédition du
%          Chap I ? ») are transcribed as \marginal.
%
% Whose hand. Only pages 48-51 and the pencil annotations on 57, 59, 60 are his.
% Pages 43-46 are Lê's, 55-60 Deligne's, 52-53 a typescript. The transcription
% keeps each writer's text as written; nothing is normalised across writers.
%
% Notation fixed for the batch. Script S, O, K of the errata become
% \mathcal{S}, \mathcal{O}_X, \mathcal{K}_X; the fraktur Z of 21.6.2 is
% \mathfrak{Z}; Deligne's blackboard N in A_s^N is \mathbf{N}; his "lim" with an
% arrow under it is \varprojlim, and the quoted pro-object limit keeps his
% quotation marks. His J-adic J and his radical R(B) are set as written.
%
% What is NOT here. No attempt was made to check the errata against the
% printed EGA IV, to identify the addressee of the typescript, or to follow the
% Deligne letter past page 60.
\documentclass[11pt,a4paper]{article}
% Le préambule commun à toutes les transcriptions du fonds Grothendieck.
%
% Il tient en une idée : **l'incertitude se compose, elle ne se résout pas.**
% Une transcription qui lisse un mot douteux a détruit la seule chose pour
% laquelle on la faisait — un lecteur doit pouvoir savoir, à chaque endroit,
% ce qui était sur le feuillet et ce qui a été ajouté. Les six macros qui
% suivent sont donc l'appareil critique, et non de la décoration.
%
% Les mêmes commandes sont reconnues par `scripts/render.mjs`, qui produit la
% vue de lecture du site. Une macro ajoutée ici doit l'être aussi là-bas,
% faute de quoi le rendu échouera bruyamment — ce qui est voulu : un rendu
% silencieusement faux est pire qu'un rendu absent.

\ProvidesPackage{grothendieck}[2026/08/08 Transcriptions du fonds Grothendieck]

\usepackage{amsmath,amssymb,amsthm}
\usepackage{mathrsfs}
\usepackage{stmaryrd}
% Les diagrammes commutatifs sont partout dans ce fonds : c'est la langue
% même de la géométrie algébrique de Grothendieck, et une transcription qui
% les rendrait en prose ne transcrirait rien.
\usepackage{tikz-cd}
% Français par défaut — c'est la langue des feuillets — mais l'anglais est
% chargé aussi : la traduction et le résumé sont des documents anglais, et la
% typographie française (espaces avant les deux-points) y serait une faute.
% Ces documents basculent par \selectlanguage{english} en tête de corps.
\usepackage[english,french]{babel}
\usepackage{fontspec}
\usepackage{geometry}
\geometry{a4paper,margin=25mm,marginparwidth=28mm}
\usepackage{marginnote}
\usepackage{xcolor}
% ulem, not soul. `soul`'s \st and \ul are built on a character-by-character
% scan that XeLaTeX's font machinery breaks: the document fails with an
% undefined control sequence rather than degrading. ulem does the same two jobs
% and is Unicode-engine safe. `normalem` stops it hijacking \emph, which the
% transcriptions use for Grothendieck's own emphasis.
\usepackage[normalem]{ulem}
\usepackage{hyperref}
\hypersetup{colorlinks=true,linkcolor=black,urlcolor=black,pdfborder={0 0 0}}

% --- Métadonnées du lot -----------------------------------------------------
% Reprises telles quelles par le script de rendu, qui les affiche en tête de
% la vue de lecture. Elles fixent ce que le fichier prétend couvrir : sans
% elles, une transcription flotte sans point d'attache dans un fonds de seize
% mille feuillets.

\newcommand{\folder}[1]{\gdef\grothFolder{#1}}
\newcommand{\batch}[1]{\gdef\grothBatch{#1}}
\newcommand{\leaves}[2]{\gdef\grothFirst{#1}\gdef\grothLast{#2}}
\newcommand{\foldertitle}[1]{\gdef\grothFoldertitle{#1}}
\gdef\grothFolder{?}\gdef\grothBatch{?}\gdef\grothFirst{?}\gdef\grothLast{?}\gdef\grothFoldertitle{}

% --- Le résumé --------------------------------------------------------------
%
% Il ouvre la lecture modernisée, sous le titre « Résumé », et il est composé
% en petit. Ce n'est pas un ornement : c'est le seul passage du document qui
% ne soit pas des mathématiques, et un lecteur doit pouvoir le reconnaître
% avant de l'avoir lu — puis le sauter s'il connaît déjà le sujet, ou s'y
% arrêter s'il ne le connaît pas. Le filet qui le suit marque la reprise du
% corps.

\newenvironment{resume}
  {\small\setlength{\parskip}{0.45em}}
  {\par\medskip\hrule\medskip}

% --- Mots-clés ---------------------------------------------------------------
%
% En anglais, en clôture du résumé de la lecture modernisée : le vocabulaire
% moderne sous lequel chercher ce que les feuillets construisent. Le manifeste
% du site (`npm run manifest`) les extrait pour étiqueter la cote entière —
% cette ligne est la source unique des tags, il n'y a pas de fichier à côté.
\newcommand{\keywords}[1]{\par\smallskip\noindent
  {\footnotesize\textsc{Keywords} --- \textit{#1}}\par}

% --- L'appareil critique ----------------------------------------------------

% Le numéro de feuillet, en marge. C'est l'ancre qui permet de régler un
% désaccord sur une formule en regardant *un* feuillet plutôt qu'une chemise
% de six cents. Le site s'en sert aussi pour tourner les pages du fac-similé
% quand on fait défiler la transcription.
\newcommand{\leaf}[1]{%
  \marginnote{\footnotesize\color{gray!70}\textsf{#1}}%
  \label{leaf:#1}\ignorespaces}

% Illisible : on ne devine pas. Un mot inventé qui se lit comme les autres est
% le pire résultat possible.
\newcommand{\ill}{\textcolor{red!60!black}{[\,\ldots\,]}}

% Lecture douteuse : le mot est donné, et signalé comme incertain.
\newcommand{\uncertain}[1]{\uline{#1}}

% Ajout de l'éditeur — une lettre manquante, un mot sous-entendu.
\newcommand{\add}[1]{\textcolor{blue!50!black}{[#1]}}

% Biffé par Grothendieck lui-même. Ce qu'il a rayé fait partie du document :
% une rature dit souvent où la pensée a changé de direction.
\newcommand{\struck}[1]{\sout{#1}}

% Note du transcripteur, sur l'état matériel du feuillet.
\newcommand{\note}[1]{\par\smallskip{\small\color{gray!60!black}\textit{[#1]}}\par\smallskip}

% Note marginale de Grothendieck, distincte du corps du texte.
\newcommand{\marginal}[1]{\par\smallskip\noindent\hspace{1em}%
  {\small\color{gray!60!black}#1}\par\smallskip}

% --- Titre ------------------------------------------------------------------

\AtBeginDocument{%
  \begin{center}
    {\large\bfseries Fonds Alexandre Grothendieck --- cote n\textsuperscript{o}~\grothFolder}\\[2pt]
    {\itshape\grothFoldertitle}\\[4pt]
    {\small feuillets \grothFirst--\grothLast\ (lot \grothBatch)}\\[2pt]
    {\footnotesize\color{gray!60!black}Transcription établie d'après le fac-similé numérisé
     par l'Université de Montpellier.\\
     Les lectures incertaines sont soulignées, les illisibles notées [\,\ldots\,],
     les ajouts entre crochets.}
  \end{center}
  \vspace{1em}}

\endinput


\folder{23}
\batch{3}
\pages{41}{60}
\dating{1966-1968}
\watermark{Édition de démonstration}
\foldertitle{EGA IV. Compléments : notes manuscrites (s.d.), lettres (1966, 1968), tapuscrit (s.d.).}

\begin{document}

\section{Lettre et errata de Lê Dũng Tráng (11 janvier 1968) : EGA IV, §§ 20--21}

\page{43}
\note{Lettre de Lê Dũng Tráng, de sa main, datée « Paris le 11/01/68 »,
adressée « à Monsieur le Professeur A. Grothendieck, IHES, Bures s/Yvette ».
Elle accompagne les trois feuillets d'errata qui suivent, numérotés ①, ②, ③
par lui en haut à gauche.}

Monsieur le Professeur,

Ayant égaré les feuilles sur lesquelles étaient relevés les errata, aussi je
me permets de vous envoyer si tardivement ces nouvelles-ci. Les errata ont été
relevés pour la plupart dans les définitions. Je ne sais pas s'ils entravent
des démonstrations car celles-ci faisaient intervenir des éléments de
mathématiques que je ne possède pas comme les idéaux associés ou la
profondeur d'un anneau. Pardonnez donc la trivialité de ces errata.

Veuillez agréer, Monsieur le professeur, l'expression de mes sentiments
respectueux et de ma profonde admiration.

\page{44}
\note{Feuillet ① des errata, de la main de Lê Dũng Tráng. La première
indication (« A la place de \ldots\ lire : ») est ajoutée en petit au-dessus
de la première ligne, après coup ; elle est rendue ici à sa place logique.}

(20.1.3) A la place de « Nous allons nous intéresser\ldots\ n'est autre que
l'ensemble des éléments réguliers de $\mathcal{O}_X$ » lire : « Nous allons
nous intéresser ici au cas où $\mathcal{S}$ est le sous-faisceau
$\mathcal{S}(\mathcal{O}_X)$ de $\mathcal{O}_X$ tel que pour tout ouvert $U$,
$\Gamma(U, \mathcal{S})$ soit l'ensemble des éléments réguliers de l'anneau
$\Gamma(U, \mathcal{O}_X)$ dont la restriction à tout ouvert $V \subset U$ est
encore régulière dans $\Gamma(V, \mathcal{O}_X)$. Il est immédiat qu'il s'agit
d'un faisceau car un élément est dans $\Gamma(U, \mathcal{S})$ si et seulement
si il est dans $\Gamma(U, \mathcal{O}_X)$ et \struck{tous} ses germes en tout
point $x$ de $U$ sont des éléments réguliers de $\mathcal{O}_{X,x}$. Il faut
remarquer que $\mathcal{S}(\mathcal{O}_X)_x$ n'est pas \struck{toujo}
nécessairement l'ensemble des éléments réguliers de $\mathcal{O}_{X,x}$. »

(20.1.4) au lieu de : « qui n'est autre que \struck{un} \add{l'}anneau total
des fractions de $\Gamma(U, \mathcal{O}_X)$ », --- lire : « qui est un anneau
de fractions à dénominateurs dans une partie multiplicativement stable
composée d'éléments réguliers ».
\note{le « l' » est de la main de Lê, porté au-dessus du « un » biffé ; il
n'est pas de l'éditeur.}

(20.1.8) à la place de : « en effet, si $s \in \Gamma(U, \mathcal{S}(\mathcal{K}_X))$
\ldots\ inversible dans cet anneau de fractions », --- lire :
« en effet, si $s \in \Gamma(U, \mathcal{S}(\mathcal{K}_X))$, pour tout
$x \in U$ il existe un voisinage ouvert $V \subset U$ de $x$ tel que $s|V$ soit
un élément régulier de $\Gamma(V, \mathcal{O}_X)[\Gamma(V, \mathcal{S})^{-1}]$
dont toutes les restrictions à des ouverts $W$ inclus dans $V$ sont des
éléments réguliers de $\Gamma(W, \mathcal{O}_X)[\Gamma(W, \mathcal{S})^{-1}]$
et on sait que dans ce cas $s|V$ est inversible dans
$\Gamma(V, \mathcal{O}_X)[\Gamma(V, \mathcal{S})^{-1}]$. »

(20.1.11) Remplacer « $\mathcal{S}_f(U)$ l'ensemble des sections régulières
$s \in \Gamma(U, \mathcal{O}_X)$ \struck{telles que} », --- par :
« $\mathcal{S}_f(U)$ l'ensemble des sections régulières
$s \in \Gamma(U, \mathcal{S})$ \struck{telles que} de $\mathcal{O}_X$ au
dessus de $U$ telles que\ldots\ »
\note{au-dessus du premier « telles que » biffé, un « que » est repris en
interligne ; il ne modifie pas la lecture.}

\page{45}
\note{Feuillet ② des errata.}

(21.1.5) A la place de : « Le faisceau $\mathcal{S}(\mathcal{O}_X)$ \ldots\
les éléments réguliers de $\Gamma(U, \mathcal{O}_X)$ », lire : « Le faisceau
$\mathcal{S}(\mathcal{O}_X)$ (20.1.3) dont les sections au-dessus d'un ouvert
$U$ de $X$ sont les éléments réguliers de $\Gamma(U, \mathcal{O}_X)$ dont
toutes les restrictions aux ouverts $V$ inclus dans $U$ sont encore des
éléments réguliers ».

(21.1.6) A la fin du paragraphe au lieu de : « il existe un élément régulier
$t$ de $\Gamma(U, \mathcal{O}_X)$\ldots\ » lire « il existe un élément
\struck{régulier} $t$ de $\Gamma(U, \mathcal{S}(\mathcal{O}_X))$ ».

(21.2.9) Après la ligne (21.2.9.1) lire :

« qui est un isomorphisme : en effet si $U$ est un ouvert de $X$ tel que
$\mathcal{J}|U = \mathcal{O}_U \cdot f$ où $f \in \Gamma(U, \mathcal{K}_X^*)$,
d'après (21.2.2) on déduit que l'on a \struck{un isomorphisme} restreint à $U$
un isomorphisme \struck{de} qui à toute section $t$ de
$\mathcal{K}_X(\mathcal{J})|U$ au-dessus de $V \subset U$ fait correspondre une
section $\bar{t}$ de $\mathcal{K}_X|U$ au-dessus de $V \subset U$ ; restreint à
$U$ l'isomorphisme (21.2.9) est celui qui à toute section $t$ de
$\mathcal{K}_X(\mathcal{J})|U$ au-dessus de $V \subset U$ fait correspondre la
section $\bar{t}\,(f|V)$ au-dessus de $V \subset U$. Dans l'isomorphisme
(21.2.9.1) \ldots\ »

(21.2.9) Après la ligne (21.2.9.2) lire :

« et on désigne par $s_D$ la section méromorphe régulière de
$\mathcal{O}_X(D)$ au-dessus de $X$ qui correspond par cet isomorphisme à la
section $1$ de $\mathcal{K}_X$. Si $U$ est un ouvert de $X$ tel que
$\mathcal{O}_X(D)|U = \mathcal{O}_U \cdot f^{-1}$ où
$f \in \Gamma(U, \mathcal{K}_X^*)$ \struck{et si} on a $D|U = \operatorname{div}_U(f)$.
De plus pour un tel $U$ on a un isomorphisme de $\mathcal{K}_X(\mathcal{J})|U$
sur $\mathcal{K}_X|U$ qui à toute section $t$ au-dessus de $V \subset U$ fait
correspondre $\bar{t}$ ; cet isomorphisme fait correspondre à $s_D|U$ la
section $f$ de $\Gamma(U, \mathcal{K}_X^*)$ donc :
\[
  (21.2.9.3) \qquad \operatorname{div}(s_D) = D .
\]
\note{dans « $\mathcal{O}_X(D)|U = \mathcal{O}_U \cdot f^{-1}$ » le « $(D)$ »
et l'exposant $-1$ sont repassés à l'encre par-dessus une première écriture ;
la lecture retenue est celle de la correction.}

\page{46}
\note{Feuillet ③ des errata.}

(21.2.12) à la première ligne ajouter \struck{l'hypothèse} \struck{condition} :
« localement noethérien » entre « préschéma » et « $X$ ».

\struck{(21.3.6)}
\note{le numéro est barré de hachures croisées et rien ne le suit : l'erratum
prévu a été abandonné.}

(21.6.2) fin du paragraphe lire $\mathfrak{Z}^p(X)$ au lieu de $Z^p(X)$.

\section{EGA IV 18 ; EGA VI}
\note{Un seul feuillet, de sa main, à l'encre, rapide ; l'intitulé de la
section reprend ses deux titres, « EGA IV 18 » en tête et « EGA VI » au milieu
du feuillet. Le second bloc est séparé du premier par un trait horizontal.}

\page{48}

\begin{tikzcd}
X \arrow[d, "u"'] \\
Y
\end{tikzcd}

non ramifié

$U$ \struck{\ill{}} ouvert schématiquement dense de $Y$

$Y_0$ sous-schéma fermé de $Y$, tel que $|Y_0| = |Y|$

$u_U$ une \struck{\ill{}} section de $X$ sur $U$

$u_{Y_0}$ une section de $X$ sur $Y_0$
\note{la seconde ligne est écrite par un trait de rappel sous « une section
de $X$ sur », suivi de « $Y_0$ » : le « une section de $X$ sur » est
sous-entendu par lui, non ajouté par nous.}

$u_U$ et $u_{Y_0}$ coïncident sur $U_0 = Y_0 \cap U$.

Alors $\exists !$ section $u$ de $X$ sur $Y$, dont les restr.\ :
$Y_0$ et $U \mapsto u_{Y_0}$ et $u_U$.

EGA VI

\begin{tikzcd}[column sep=small]
X \arrow[d, "\text{non ram.}"'] & & \\
Y & Y' \arrow[l, "\text{loc. noeth. fini épim.}"'] & Y'' \arrow[l, bend left=20] \arrow[l, bend right=20]
\end{tikzcd}

$Y'' \to Y' \times_Y Y'$ surjectif
\note{le « noeth. » sous la flèche $Y' \to Y$ est surchargé et se lit comme
biffé --- ce serait cohérent avec la note marginale qui suit --- mais la
biffure n'est pas certaine et le mot est laissé dans le diagramme.}

\marginal{On doit pouvoir éliminer l'hyp.\ loc.\ noeth.\ et remplacer fini
par \uncertain{entier}}
\note{note verticale en marge droite, à hauteur du diagramme, le dernier
mot souligné.}

Alors $X(Y) \to X(Y') \rightrightarrows X(Y'')$ est exact

\emph{Corollaire des deux} \quad Variante en remplaçant $Y_0$ par un $Y'$
fini surjectif sur $Y$

\emph{Corollaire} \quad Si $f : X \to Y$ est \ill{} \ill{}\ldots
\note{la dernière ligne s'interrompt sur des points de suspension de sa
main ; les deux mots entre « est » et les points n'ont pas pu être lus (le
premier pourrait être « un. », le second commence par « f »).}

\section{Suites régulières --- Th.\ 11.2.9 de Raynaud}
\note{L'intitulé est de sa main, sur un feuillet par ailleurs blanc (la page
49), en haut à droite, « Suites régulières » souligné : c'est la couverture
du dossier qui suit.}

\page{50}
\note{Feuillet à en-tête de l'IHES, de sa main, au crayon.}

\emph{Applications de 11.2.9 dans la suite du Chap.\ IV}

11.3.4. (i)

11.3.8. b')

17.12.1. c')

19.2.4. c')

19.4.6. (ii)

À signaler \struck{\ill{}} remarques ?
\note{« À signaler » et « remarques » sont soulignés ; le mot biffé entre
eux est court et ne se lit pas.}

\page{51}
\note{Feuillet à en-tête de l'IHES, de sa main, à l'encre.}

\emph{Suites régulières, immersions régulières}

$0_{\mathrm{IV}}$ 15

IV 11.3.7. \quad 11.3.8.

IV 16.9.

IV 19.

IV 21.15.

Lettres Deligne 17.\ Oct.\ 66.\ et 29 Oct.\ 66.

\struck{Notion d'immersion régulière dans le cas non noethérien :
\uncertain{est-elle vraisemblable} ?? Vérifier § 16 !!}
\note{ce bloc de trois lignes est encadré de deux traits horizontaux et
barré de deux longues diagonales ; il est transcrit comme biffé en entier.
Les lettres Deligne des 17 et 29 octobre 1966 sont celles des pages 57--60.}

\page{52}
\note{Tapuscrit, copie carbone, deux pages (52--53), sans en-tête ni
signature ; le texte est à la première personne (« mon séminaire sur
Riemann-Roch ») et se lit comme un post-scriptum de lettre. Les surfrappes
de la machine deviennent \struck{}, les mots frappés en interligne \add{}.}

PS \quad Dans mon séminaire sur Riemann-Roch cette année, on doit
mani\struck{f}puler systématiquement les immersions régulières dans le cas
non noethérien, et il apparait (hélas) que la notion introduite sous ce nom
dans EGA 16.9.2.\ n'est pas la bonne notion. Un premier canular important,
c'est que si $A$ est un anneau local, et $x, y$ deux éléments de son idéal
maximal, il est facile de donner un exemple ou $(x, y)$ est une suite
régulière, et $(y, x)$ ne l'est pas. La notion de suite régulière de
générateurs d'un idéal ne jouit donc pas du tout des propriétés habituelles,
comme celle qu'il suffise de prendre une base de $J/J^2$ et de relever. Par
suite, la notion d'immersion régulière n'a pas l'air de devoir se comporter
comme il faut par localisation plate ou même étale, ce qui est un autre grave
inconvénient, en entrainant toute sorte d'autres d'ailleurs. Il apparait
maintenant que la notion vraiment serviable pour une suite $(f_1, \ldots, f_n)$
d'éléments d'un anneau (engendrant un idéal $J$) c'est que le complexe de
l'algèbre extérieure construit avec soit nul en degrés $0$. Alors on prouve
que la suite est aussi quasi-régulière, et de plus que toute autre suite de
\add{n} générateurs de $J$ a la même propriété que celle énoncée pour
$(f_1, \ldots, f_n)$. Il y aurait lieu alors de dire que $J$ est un idéal
régulier si localement il peut être engendré par une suite comme dessus
(qu'on pourrait peut-être appeler « suite Koszul-régulière » ou
« essentiellement régulière », si on tient à ne pas changer la terminologie
des \emph{suites} régulières), et à définir de façon correspondante la
notion d'immersion régulière. Ces notions (suite K-régulière, idéal régulier,
immersion régulière) se descendent bien par fpqc. Un composé de deux
immersions régulières est itou, et les deux autres propriétés de
transitivité de 19.1.5.\ (iii) revu semblent marcher encore sans hypothèse
noethérienne (Berthelot en a déjà vérifié une,
\note{« complexe de l'algèbre extérieure construit avec soit nul » : le
« avec » est suivi directement de « soit », sans le nom du complexe ; frappé
ainsi. Le « n » de « suite de $n$ générateurs » est frappé en interligne.}

\page{53}

en tous cas. De plus, bien entendu, pour ce qu'on veut en faire pour R.R., c'est
bien la nouvelle notion d'immersion régulière qui fait tout marcher.

Du point de vue pratique, il me semble évidemment trop tard pour faire les
changements désirables sur épreuves dans IV 16, 19. On peut cependant ajouter
une remarque 19.1.8., dans laquelle nous ferions l'aveu que \struck{la} la
définition adoptée dans le texte de immersion régulière \struck{n}s'est
\struck{pas} avérée peu maniable dans le cas non noethérien, et qu'il y aurait
lieu de la modifier comme je viens de le dire, en référant le lecteur
intéressé, \add{pour des détails} en attendant la réédition du Chap IV des
EGA, au séminaire SGA 6 (qui devrait sortir vers la fin de l'année 67 ou
début 68). De plus, à l'endroit du par.\ 16 ou on introduit la définition
incriminée, nous pourrions insérer une note de bas de page, renvoyant le
lecteur à la remarque 19.1.8.\ concernant cette définition. (NB Dans le
par.\ 16, c'est en fait la notion quasi-régulière qui sert vraiment \ldots).
\note{le passage « ajouter une remarque 19.1.8. \ldots\ la définition
incriminée » est marqué en marge gauche d'un double trait vertical au
crayon.}

Au moment de la réédition du Chap IV, il y aura lieu également de revoir
$0_{\mathrm{IV}}$ 15 dans le même esprit ; je prévois que l'on pourra
également y remplacer partout les hypothèses « suite régulière », par
l'hypothèse plus faible et plus maniable « suite Koszul-régulière ». S'il en
est bien ainsi, il s'imposera alors, pour les suites également, de réserver
le mot le plus court à la notion qui manifestement est la plus serviable, et
il faudra \uncertain{envisaher}\note{sic} de changer \struck{\ill{}} notre
terminologie, en disant « suite régulière » au lieu de « suite
Koszul-régulière », et en réservant tout au plus un terme tel que « suite
strictement régulière » à la notion éculée. Bien sûr, dans le cas
noethérien, les trois variantes (régulière, quasi, et Koko) sont
équivalentes.

\section{Lettres Deligne}
\note{L'intitulé est de sa main, en haut à droite d'un feuillet par
ailleurs blanc (la page 54) : couverture des lettres qui suivent. Les trois
lettres sont de la main de Pierre Deligne, à l'encre, et sont classées dans
l'ordre inverse des dates ; les annotations au crayon qu'elles portent sont
de Grothendieck et sont rendues par \marginal{}.}

\page{55}

Bruxelles 5, le 15 décembre 1966

Cher Grothendieck,

Je lis aujourd'hui ta lettre du 28. Après avoir tourné 7 fois ma langue dans
ma bouche, je répète mon objection à l'hypothèse « tout quotient de tout
sous-module de $M$ est séparé pour la topologie $J$-adique », qui
\emph{n'est pas} \add{en général} vérifiée pour $M$ de type fini sur la
$A$-algèbre $B$, $A$ et $B$ noethériens, $JB \subset R(B)$.
\note{« en général » est ajouté par Deligne en interligne, avec un signe
d'insertion ; le \add{} le signale comme ajout de l'auteur de la lettre,
non de l'éditeur. « n'est pas » est souligné.}

\emph{Contre exemple} : $A$ noethérien local, $J = \mathfrak{m}$

$B = A[[T]]$, $M = B$

comme \emph{$A$-module}, $B \simeq A_s^{\mathbf{N}}$, donc tout module
quotient de $A_s^{(\mathbf{N})}$ est quotient d'un sous-$A$-module de $B$.
\struck{Un} Un module ayant un système dénombrable de générateurs n'a aucune
raison d'être séparé !
\note{l'indice $s$ de $A_s^{\mathbf{N}}$ et $A_s^{(\mathbf{N})}$ est écrit
ainsi et n'est pas expliqué dans la lettre.}

Bien sûr, les démonstrations \struck{ne} n'utilisent que des quotient et
sous-modules comme $B$-modules.

J'essayerai encore une fois de lire 19.7.

Bien à toi

P. Deligne

P.S. Je rappelle que l'adresse où tu me touches le plus rapidement est :
S.M. P. Deligne, Athénée Royal de Rösrath, BPS 8, FBA (forces belges en
Allemagne)

P.S. Je commence à lire IV § 16. La démonstration de 16.1.5 est canulée,
comme on le voit en prenant $X = \operatorname{Spec}(k[T])$,
$Y = \operatorname{Spec}(k[T]) \sqcup \operatorname{Spec}(k)$, tous deux
affins.
\note{le « 5 » de 16.1.5 est repassé en gras par-dessus un premier chiffre.
En marge gauche, un croquis : deux accolades « $Y$ » et « $X$ » devant deux
traits horizontaux superposés, celui de $Y$ portant un point isolé au-dessus,
et une flèche verticale descendante de $Y$ vers $X$ --- la droite affine
avec un point disjoint, au-dessus de la droite affine.}
Avec les notations de 16.1.15 (i), $J = 0$, mais les voisinages
infinitésimaux de $Y$ ne sont pas \struck{\ill{}} triviaux en
$\operatorname{Spec}(k)$

J'ignore si la proposition est vraie ---

\page{56}

Par ailleurs, même en admettant 16.1.5, 16.1.6 n'est pas convaincant en
général : si $B$ est de présentation finie sur $A$, j'ignore si le noyau de
$A \to B$ est un idéal de type fini. C'est sans doute faux.

\page{57}

Rösrath, le 29 octobre 1966

Cher Grothendieck,

Je te signale tout d'abord mon adresse en Allemagne, maintenant définitive
(jusqu'au 1er juillet 67) : S.M. Pierre Deligne, Athénée
royal, BPS 8, FBA (forces belges en Allemagne).

Je suis en train de lire EGA IV 19 que tu m'as envoyé.

\textbf{a) remarque sur les hypothèses de séparation utilisées}

L'hypothèse $(*)$ « tout quotient de tout sous-module de $M$
\struck{\ill{}} (resp $M^n$) est séparé pour la topologie $J$-adique » me
semble idiote :

(i) elle \struck{n'implique pas} n'est pas impliquée par l'hypothèse de
finitude relative noethérienne standart : « $M$ fini sur la $A$-algèbre
noethérienne $B$, $A$ noethérien, et $JB \subset R(B)$ » qui elle est
naturelle.
\marginal{?}
\note{un point d'interrogation au crayon, en marge gauche, contre un trait
vertical qui couvre l'énoncé (i).}

(ii) le seul exemple non trivial indiqué où cette hypothèse est vérifiée,
soit 19.5.6 (ii), est canulé, car \struck{\ill{}} dans $(*)$ on ne se limite
pas à prendre des quotients par des sous-modules gradués.
\marginal{Err}
\marginal{\ill{}}
\note{« Err », souligné, au crayon, en marge gauche à hauteur de (ii) ;
au-dessous, deux mots au crayon écrits en oblique et soulignés d'un trait
montant, dont le premier pourrait se lire « ja » et le second n'a pas été
lu ; une parenthèse ouvrante couvre les deux dernières lignes de (ii).}

[exemple : $A = B[T]$, $M = N[T]$, $N$ est un $A$-module, $T$ agissant par
$1$, et ce $A$-module est quotient de $M$]
\marginal{!}
\note{un point d'exclamation au crayon en marge gauche de l'exemple ; le
$N$ de « $N$ est un $A$-module » est repassé sur une autre lettre.}

Je crois qu'il faut éviter d'employer des hypothèses de séparation qu'on ne
rencontre pas dans la nature, comme « idéalement séparé » etc., et se
limiter à utiliser (i) --- J'avoue toutefois, avec honte, que dans le même
geste sacrilège je serais prêt à supposer tous les schémas séparés et
noethériens\ldots

Sinon, pour coiffer (i) et $(*)_{\mathrm{resp}}$ on pourrait prendre :

$\forall f : N \to P$, si $N$ et $P$ sont finis sur $A$,
$\operatorname{Ker}(M \otimes N \to M \otimes P)$ est $J$-séparé.
\marginal{!}
\note{un point d'exclamation au crayon et une parenthèse en marge gauche
contre l'énoncé. Le $P$ y est écrit trois fois par-dessus un $M$ ; dans le
noyau, ce qui précède $M \otimes P$ est biffé et surchargé, et l'on retient
la lecture corrigée.}

\page{58}

\textbf{b) réponse à 19.5.6.\ (iii)}

l'implication b) $\Rightarrow$ c) de 19.5.5 est vraie sans qu'on doive
faire aucune hypothèse de séparation. Pour le voir, il suffit de lire le
cas $n = 1$ de la démonstration « à la Schwartz ».

Vu l'urgence que tu m'as signalée, et bien qu'il soit peut-être trop tard,
je joins une rédaction de la démonstration (copiée sur la primitive)

--- Je n'ai pas encore compris 19.7.

--- J'ai beaucoup aimé le projet EGA I § 0 jusqu'à la pg 10, où la topologie
du spectre tombe du ciel. Ne faudrait-il pas parler de topologie avant de
parler d'espace topologique ?

--- Je ne suis pas capable de rédiger une théorie naïve utilisable des
intersections, et n'ai pas d'autre remarque à faire sur le « plan des EGA »
pour le moment

Bien à toi

P. Deligne

\page{59}
\note{Lettre paginée 1 et 2 par Deligne (pages 59 et 60) ; elle s'interrompt
au bas de la page 60 au milieu d'une démonstration et se poursuit dans le
lot suivant.}

Bruxelles 5, le 17 octobre 1966

Cher Grothendieck,

J'ai bien reçu ta lettre.

Je n'ai pas l'intention d'essayer de démontrer, cette année, le « théorème
de Lefschetz vache ». J'espère donc que tu trouves quelqu'un qui le fasse.

J'ai été malade au début du mois et, comme lecture de convalescence, je lis
EGA IV. Je ne sais pas si tu collectionnes déjà les errata et compléments
pour EGA $0_{\mathrm{IV}}$ ; je t'envoie toujours ceux qui suivent.

\textbf{$0_{\mathrm{IV}}$ 19.3.12} \quad est \emph{canulé} : il faut prendre
pour $W_\lambda$ l'ensemble des homomorphismes d'\emph{algèbre}
$B \to C_\lambda$ qui relèvent $B \to C_\lambda / J_\lambda$, et ce n'est pas
une variété linéaire [on n'a pas $J_\lambda^2 = 0$] --- Pour faire marcher
l'argument, il faudrait savoir que pour $A$ artinien, l'ensemble des parties
de $A^I$ définies par \struck{\ill{}} des équations $P(x_i) \in \mathfrak{b}$
($\mathfrak{b}$ idéal) jouit d'une propriété de compacité. C'est faux pour
$I$ infini, déjà pour $A$ un corps infini [$I = k \cup \{\infty\}$, parties
$(x_\infty - a)\, y_a = 1$, les $\cap$ finies sont non vides, l'$\cap$ est
vide] --- Je crois qu'il n'y a pas d'espoir.
\marginal{oui}
\note{« oui » au crayon, entouré, en marge gauche sous le numéro
$0_{\mathrm{IV}}$ 19.3.12 ; « canulé » et « algèbre » sont soulignés par
Deligne.}

\textbf{$0_{\mathrm{IV}}$ 17.3.4} \quad Je te signale une jolie démonstration
cohomologique de ce résultat, qui fournit un résultat plus complet (lequel se
trouve dans Northcott)

\marginal{Th.\ de Auslander}
\note{au crayon, entre parenthèses, en marge gauche devant « prop ».}

\emph{prop} : $A$ \emph{local noethérien}, $M$ de type fini sur $A$, non
nul, et de \emph{dimension projective finie}. Alors
\[
  \dim\operatorname{proj}(M) = \operatorname{prof} A - \operatorname{prof} M .
\]

on utilise la formule suivante, valable pour un complexe parfait $L$
\[
  R\operatorname{Hom}(K, L) = R\operatorname{Hom}(K, A) \overset{L}{\otimes} L
\]
avec $K = k$, $L = M$. On trouve une suite spectrale
\[
  E_2^{pq} = \operatorname{Tor}_{-p}\big(\struck{\ill{}}\operatorname{Ext}^q(k, A), M\big)
  \Rightarrow \operatorname{Ext}^{p+q}(k, M) . \qquad \text{(dessin ci-contre)}
\]
si $d = \dim\operatorname{proj} M$, $a = \operatorname{prof} A$, on trouve
\[
  \operatorname{Ext}^n(k, M) = 0 \text{ si } n < a - d \text{ et}
\]
\[
  \operatorname{Ext}^{a-d}(k, M) = \operatorname{Tor}_d\big(\operatorname{Ext}^a(k, A), M\big),
  \text{ non nul}
\]
par hypothèse (car $\operatorname{Ext}^a(k, A)$ est somme de copies de $k$)

D'où la profondeur de $M$\ldots
\note{le dessin « ci-contre », en bas à gauche : un axe horizontal $p$
(flèche vers la droite) et un axe vertical $q$ ; une bande verticale hachurée
de largeur $d$, mesurée à gauche de l'axe des $q$, commence à la hauteur $a$
au-dessus de l'axe des $p$ et monte --- le support de $E_2^{pq}$, nul pour
$q < a$ et pour $p < -d$.}

\page{60}
\note{page 2 de la lettre, paginée ainsi par Deligne.}

Le même résultat, \add{en remplaçant $\dim\operatorname{proj}$ par
$\operatorname{Tor}\dim$}, est encore valable pour $M$ de type fini sur un
anneau \struck{local} noethérien $B$, qui soit une $A$-algèbre telle que
$\mathfrak{m} B \subset \operatorname{radical}(B)$. Il faut savoir (i) qu'on
peut encore définir la profondeur de $M$ sur $A$ dans ce cas (ii) que la
$\operatorname{Tor}$ dimension de $M$ est le plus grand entier $d$ tel que
$\operatorname{Tor}^d(k, M) \neq 0$, où $k$ est le corps résiduel de $A$.
\note{« en remplaçant dim proj par Tor dim » est ajouté par Deligne en
interligne, avec un signe d'insertion ; le $\operatorname{Tor}^d$ est écrit
avec l'indice en exposant, comme ici.}

C'est là une conséquence facile du résultat suivant, pour lequel je ne
connais pas de référence :

\emph{prop} : Soit $A$ un anneau \emph{local noethérien}, $B$ une
$A$-algèbre locale noethérienne ($A \to B$ étant local),
\emph{$\mathfrak{m}$ un idéal de $A$}, \struck{$M$}$N$ un $A$-module de
\emph{type fini}, $M$ un \struck{$A$}$B$-module de type fini. On a
\[
  \varprojlim \operatorname{Tor}_p^A(N / \mathfrak{m}^k N, M)
  = \operatorname{Tor}_p^A(N, M)^{\wedge}
\]
et plus précisément une égalité de pro-objets
\[
  \underset{\longleftarrow}{\text{“lim”}}\;
  \operatorname{Tor}_p^A(N, M) \big/ \mathfrak{m}^k \operatorname{Tor}_p^A(N, M)
  \;\xrightarrow{\ \sim\ }\;
  \underset{\longleftarrow}{\text{“lim”}}\;
  \operatorname{Tor}_p^A(N / \mathfrak{m}^k N, M)
\]
\note{Deligne écrit « lim » entre guillemets, la flèche sous le mot, pour
la limite projective prise comme pro-objet.}
\note{les lettres $N$ et $B$ de l'énoncé sont écrites par-dessus un $M$ et
un $A$ : la correction est de Deligne.}

Si $B$ était plat sur $A$, on prendrait une résolution de $M$ sur $B$, et il
suffirait d'appliquer Artin-Rees, qui exprime que « le pro-objet complétion »
est foncteur exact. On peut par ailleurs remplacer $B$ par une algèbre dont
il soit quotient. On peut aussi compléter $B$. L'assertion résulte alors du
lemme suivant, qui coiffe $0_{\mathrm{IV}}$ 19.7.1.3 et $0_{\mathrm{III}}$
10.3.1
\note{les deux numéros sont portés par Deligne en marge gauche, l'un sous
l'autre, joints par « et », à hauteur de « prop » ; ils sont rapportés ici
à la fin de la phrase qu'ils complètent.}

\marginal{dans réédition du Chap I ?}
\note{au crayon, entouré, en oblique dans la marge gauche à hauteur des
conditions (i)--(iii).}

\emph{prop} : Soit $A \to B$ un homomorphisme local d'\emph{anneaux locaux
noethériens}, avec \emph{$B$ complet}. $B$ est quotient d'un anneau $C$, tel
que

(i) $C$ est local, $A \to C$ et et $C \to B$ sont locaux

(ii) \emph{$C$ est plat sur $A$}

(iii) Si $k$ est le corps résiduel de $A$, $C_0 = C \otimes_A k$,
$B_0 = B \otimes_A k$, \emph{$C_0$ est un anneau \add{régulier}}
\struck{de séries formelles sur une extension de $k$}, et l'application
graduée associée à $C_0 \to B_0$ : $\operatorname{gr} C_0 \to \operatorname{gr} B_0$
est bijective en degré $0$ et $1$
\note{« régulier » est écrit par Deligne au-dessus du passage biffé, dont
la lecture « de séries formelles sur une extension de $k$ » est celle des
mots encore visibles sous la rature. « et et » est ainsi dans (i).}

\emph{Démonstration} : Soit $\mathfrak{m}$ ($\mathfrak{n}$) l'idéal maximal
de $A$ ($B$), $K$ le corps résiduel de $B$. Soit $t_i$ une base d'un
supplémentaire de l'image de $\mathfrak{m}/\mathfrak{m}^2 \otimes_k K \to
\mathfrak{n}/\mathfrak{n}^2$. Si on relève les $t_i$ dans $B$, on définit une
flèche $A[[T_i]] \to B$. Soit $A' = A[[T_i]]$, $\mathfrak{m}'$ son idéal
maximal. $A'$ est plat sur $A$, et on a \struck{\ill{}}
\[
  \mathfrak{m}' / (\mathfrak{m}'^2 + \mathfrak{m} A')
  \;\xrightarrow{\ \sim\ }\;
  \mathfrak{n} / (\mathfrak{n}^2 + \mathfrak{m} B) .
\]
\note{la lettre s'interrompt ici, au bas du feuillet ; la suite est dans le
lot suivant (page 61 et au-delà).}

\end{document}
